% AY2012 制御工学 第2問（転記）
% 出典: 03_制御工学/original/03_制御工学/2012/2012se.pdf
% 要約: 一自由度振動系のブロック線図（1/M, 1/s, 1/s と D（速度帰還）,
%       K（位置帰還））。伝達関数・固有角振動数・減衰係数・
%       インパルス応答・単位ステップ応答を問う。

\section*{第2問}

一自由度振動系のブロック線図は下記のように表現できる. 次の問いにそれぞれ答えよ.

\begin{center}
\begin{tikzpicture}[>=Latex]
  \node (in)   at (0,0)   {$U(s)$};
  \node (sum1) [sum] at (1.3,0) {};
  \node (sum2) [sum] at (2.7,0) {};
  \node (Minv) [block] at (4.2,0) {$\dfrac{1}{M}$};
  \node (int1) [block] at (5.9,0) {$\dfrac{1}{s}$};
  \coordinate (vel) at (7.0,0);
  \node (int2) [block] at (8.1,0) {$\dfrac{1}{s}$};
  \coordinate (tap) at (9.3,0);
  \node (out)  at (10.2,0) {$Y(s)$};
  \node (D)    [block] at (5.2,-1.5) {$D$};
  \node (K)    [block] at (5.2,-3.0) {$K$};

  \draw[->] (in)   -- (sum1);
  \draw[->] (sum1) -- (sum2);
  \draw[->] (sum2) -- (Minv);
  \draw[->] (Minv) -- (int1);
  \draw[->] (int1) -- (int2);
  \draw[->] (int2) -- (out);
  \fill (vel) circle (1.4pt);
  \fill (tap) circle (1.4pt);

  % 速度帰還 D: vel → D → sum2（負帰還）
  \draw (vel) -- (7.0,-1.5) -- (D.east);
  \draw[->] (D.west) -| (sum2);

  % 位置帰還 K: 出力 → K → sum1（負帰還）
  \draw (tap) -- (9.3,-3.0) -- (K.east);
  \draw[->] (K.west) -| (sum1);

  % 符号
  \node at ($(sum1.north west)+(0.02,0.10)$) {\scriptsize $+$};
  \node at ($(sum1.south west)+(0.02,-0.10)$) {\scriptsize $-$};
  \node at ($(sum2.north west)+(0.02,0.10)$) {\scriptsize $+$};
  \node at ($(sum2.south west)+(0.02,-0.10)$) {\scriptsize $-$};
\end{tikzpicture}
\end{center}

\begin{enumerate}[label=(\arabic*)]
  \item 入力を $U(s)$, 出力を $Y(s)$ として, システムの伝達関数を求めよ.

  \item システムの固有角振動数 $\omega_n$ ならびに減衰係数 $\zeta$ を求めよ.

  \item 各パラメータを $M=1$, $D=3$, $K=2$ としたとき,
        このシステムについてのインパルス応答を求めよ.

  \item (3) のシステムについての単位ステップ応答を求めよ.
\end{enumerate}
